% !PRINTED none
% Abschnitt-opening page: no running page number visible in any of the three bands
% (top band is blank page-top whitespace; the title block sits lower on the page).
\begin{center}
\gesperrt{Fourteenth Chapter.}

\textbf{Composition of Forms and Ideals.}

\rule{0.2\linewidth}{0.4pt}

\S\,101. \gesperrt{Composition in the Orders.}
\end{center}

Let $\Delta$ be the field discriminant and $D = Q^2\Delta$ the discriminant of the order $[Q]$.

If $T = (a, b, c)$ is a primitive quadratic form of the discriminant $D$ % [sic] German source prints plural Diskriminanten.
 with $a$ positive and coprime to $Q$, then from it we can determine a basis form $\lambda$ in $O$ of an ideal $\mathfrak{a}$ coprime to $Q$, in the following way:

We set
\weq{1}{\lambda = a t_1 + \frac{b + \sqrt{D}}{2} t_2,}
\weq{2}{N(\lambda) = a(a t_1^2 + b t_1 t_2 + c t_2^2).}
If we denote the ideal determined by the functional $\lambda$, that is, the greatest common divisor of
\weq{3}{a, \frac{b + \sqrt{D}}{2}}
by $\mathfrak{a}$, then from (2) it follows that
\weq{4}{N(\mathfrak{a}) = a.}
Here $a$ is the least natural number divisible by $\mathfrak{a}$. For if this least number is $a_1$, then $a_1$ must be a divisor of $a$, say $a = a_1 a_2$. Here $a_1, a_2, b$ can have no common divisor; for, in the first place, they cannot all three be even, since otherwise we should have $a = 4a'$, $b = 2b'$, $c = c'$,
\[
D = 4(b'^2 - 4a'c)
\]
and 4 would have to divide $Q$, contrary to the assumption that $a$ is coprime to $Q$; and likewise an odd common prime divisor of $a_1, b, a_2$ would divide $Q$.
